\label{sec:task}


We study generating Wikipedia-like articles from scratch, placing emphasis on the \textit{pre-writing} stage~\citep{rohman1965pre}, which involves the demanding sub-tasks of gathering and curating relevant information (``research''). %
This models the human writing approach which has prompted
some educators to view %
Wikipedia article writing 
as an educational exercise for academic training%
~\citep{tardy2010writing}.

\begin{table}
\centering
\resizebox{\columnwidth}{!}{%
\begin{tabular}{lcccc} 
\toprule
     & \textbf{Domain} & \textbf{Scope} & \begin{tabular}[c]{@{}c@{}}\textbf{Given}\\\textbf{Outline?}\end{tabular} & \begin{tabular}[c]{@{}c@{}}\textbf{Given}\\\textbf{Refs?}\end{tabular}  \\ 
\midrule
  \citet{balepur-etal-2023-expository}   & One             & One para.      & /                                                                         & Yes                                                                     \\
 \citet{qian2023webbrain}    & All             & One para.      & /                                                                         & No                                                                      \\
  \citet{fan-gardent-2022-generating}   & One             & Full article   & Yes                                                                       & No                                                                      \\
 \citet{j.2018generating}    & All             & One para.      & /                                                                         & Yes                                                                     \\
 \citet{sauper-barzilay-2009-automatically}    & Two             & Full article   & No                                                                        & No                                                                      \\ 
\midrule
Ours & All             & Full article   & No                                                                        & No                                                                      \\
\bottomrule
\end{tabular}
}
\caption{Comparison of different Wikipedia generation setups in existing literature. Generating one paragraph does not need an article outline.}
\vspace{-0.5em}
\label{table:background}
\end{table}

\reftab{table:background} compares our work against prior benchmarks for Wikipedia generation. Existing work has generally focused on evaluating the generation of shorter snippets (\eg, one paragraph), within a narrower scope (\eg, a specific domain or two), or when an explicit outline or reference documents are supplied. A notable example is WikiSum~\citep{j.2018generating}, which treats generating Wikipedia articles as a multi-document summarization problem, with respect to the reference documents. 

Our setup emphasizes the capability of long-form grounded writing systems to research and curate content. Specifically, given a topic $t$, the task is to find a set of references $\mathcal{R}$ and generate a full-length article $\mathcal{S}=s_1s_2...s_n$, where each sentence $s_i$ cites a list of documents in $\mathcal{R}$.\footnote{In practice, $\mathcal{S}$ also includes organizational elements such as section and subsection titles, which do not require citations.}






\subsection{The FreshWiki Dataset}
\label{sec:task_dataset}

Creating a new Wikipedia-like article demands not only fluent writing but also good research skills. As modern LLMs are generally trained on Wikipedia text, we mitigate data leakage by explicitly seeking out \textit{recent} Wikipedia articles that were created (or very heavily edited) after the training cutoff of the LLMs we test. Our process can be repeated at future dates when new LLMs emerge.

To apply our date criteria, we focus on the top 100 most-edited pages, based on edit counts, for each month from February 2022 to September 2023\footnote{Obtained from \url{\detokenize{https://wikimedia.org/api/rest_v1/metrics/edited-pages/top-by-edits/en.wikipedia/all-editor-types/content/}}\{year\}/\{month\}/all-days}. 
To ensure high-quality references, we filter these articles to keep only those having B-class quality or above assessed by ORES\footnote{\url{https://www.mediawiki.org/wiki/ORES}}. We also exclude list articles\footnote{\url{https://en.wikipedia.org/wiki/Wikipedia:Stand-alone_lists}} and articles that have no subsections. While high-quality Wikipedia articles usually contain structured data (\eg, tables) and are multi-modal, we only consider the plain text component in constructing the dataset to simplify our task. More details of the dataset are in Appendix~\ref{appendix:dataset}.


\subsection{Outline Creation and Evaluation}
\label{sec:task_outline}


A full-length article is hard to generate or evaluate~\citep{xu-etal-2023-critical, krishna-etal-2023-longeval}. When human educators teach students academic writing, they sometimes supervise students at the outline stage~\citep{eriksson2015supervision} because an extensive outline indicates a comprehensive understanding of the topic and provides a solid foundation for writing the full-length article~\citep{dietz2019trec}. Inspired by this, we decompose the generation of $\mathcal{S}$ into two stages. In the pre-writing stage, we require the system to create an outline $\mathcal{O}$, which is defined as a list of multi-level section headings\footnote{Since language models process and produce sequences, we can linearize $\mathcal{O}$ by adding ``\#'' to indicate section titles, ``\#\#'' to indicate subsection titles, etc.}. In the writing stage, the system uses the topic $t$, the references $\mathcal{R}$, and an outline $\mathcal{O}$ to produce the full-length article $\mathcal{S}$.


To evaluate the outline coverage, we introduce two metrics: \textit{heading soft recall} and \textit{heading entity recall}. These metrics compare the multi-level section headings of the human-written article, considered as ground truth, and those in $\mathcal{O}$. Recognizing that an exact match between elements in these two sets of headings is unnecessary, we calculate the \textit{heading soft recall}~\citep{franti2023soft} using cosine similarity derived from Sentence-BERT~\citep{reimers-gurevych-2019-sentence} embeddings of the headings (details in Appendix~\ref{appendix:soft_heading_recall}).
We also compute the \textit{heading entity recall} which is quantified as the percentage of named entities in human-written article headings covered by $\mathcal{O}$. We extract entities with FLAIR named entity recognition (NER)~\citep{akbik-etal-2019-flair}.

